\section{Underwater Dynamics Model and Motivation}

\subsection{Classical Dynamics Background}

A standard underwater vehicle dynamics model may be written as
\begin{equation}
\bm{M}(\bnu)\dot{\bnu} + \bm{C}(\bnu)\bnu + \bm{D}(\bnu)\bnu + \bg(\boldsymbol{\eta}) = \btau,
\end{equation}
where $\bnu$ denotes body-frame velocity variables, $\boldsymbol{\eta}$ denotes pose variables, and $\btau$ denotes generalized actuation.

This form is physically meaningful but difficult to identify precisely in practice because:
\begin{itemize}
  \item thruster dynamics are nonlinear and history-dependent,
  \item hydrodynamic effects vary with operating condition,
  \item sensor observations are multi-rate and partially missing,
  \item the final use case is short-horizon prediction for control rather than full parameter recovery.
\end{itemize}

\subsection{Why a Data-Driven Surrogate}

Instead of identifying all coefficients explicitly, the current project adopts a control-oriented surrogate:
\begin{equation}
\by_{t+1:t+H} \approx f_{\theta}(\bX_t).
\end{equation}

The model is therefore asked to approximate the locally relevant short-horizon response map from aligned history windows to future observable state proxies.

\subsection{Structured Priors}

The current architecture keeps room for weak physical priors:
\begin{itemize}
  \item actuator-side lag and saturation modeling,
  \item latent hydrodynamic memory,
  \item damping-related output bias,
  \item heteroscedastic predictive uncertainty.
\end{itemize}

These priors do not replace the classical dynamics equation; they provide an intermediate layer between pure black-box prediction and full white-box identification.
